\label{chap.5}

The second-generation tester (v2) follows from the conclusions of Chapter~\ref{chap.4}. It
makes the page state explicit instead of hoping a network learns it; it remembers reached
states and returns to them exactly, so depth does not depend on luck; it deviates from the
happy path deliberately, with operators from standard test design; and it keeps learning
light and online, inside the application under test. Figure~\ref{fig:hybrid} shows the
loop. Sections~\ref{sec:abstraction}--\ref{sec:decision} describe its parts,
Section~\ref{sec:v3-options} the options added for the final tuning round, and
Section~\ref{sec:devrounds} the development rounds that shaped it.

\begin{figure}[H]
\centering
\begin{tikzpicture}[
  box/.style={draw, rounded corners=2pt, align=center, minimum height=0.9cm,
              text width=2.6cm, font=\small, execute at begin node={\hyphenpenalty=10000}},
  arr/.style={-{Latex[length=2mm]}, thick}]
  \node[box] (obs) {Observe page};
  \node[box, right=0.6cm of obs] (abs) {Abstract into a cell};
  \node[box, right=0.6cm of abs] (arc) {Archive of cells};
  \node[box, right=0.6cm of arc, fill=gray!12] (dec) {Decide\\\footnotesize bandit / curiosity / none};
  \node[box, below=1.1cm of dec] (ret) {Return to cell by exact replay};
  \node[box, left=0.6cm of ret] (op) {Apply a test operator, follow through};
  \node[box, left=0.6cm of op] (chk) {Check\\\footnotesize oracle, invariants};
  \node[box, left=0.6cm of chk] (ver) {Verify and report};
  \draw[arr] (obs) -- (abs);
  \draw[arr] (abs) -- (arc);
  \draw[arr] (arc) -- (dec);
  \draw[arr] (dec) -- (ret);
  \draw[arr] (ret) -- (op);
  \draw[arr] (op) -- (chk);
  \draw[arr] (chk) -- (ver);
  \draw[arr, dashed] (op.north) -- ++(0,0.45) -| node[pos=0.25, above, font=\footnotesize]{new pages} (obs.south);
\end{tikzpicture}
\caption{The hybrid tester's loop. Every replayed action counts towards the budget.}
\label{fig:hybrid}
\end{figure}

\section{Page abstraction}
\label{sec:abstraction}

Each page becomes a small symbolic \emph{cell}, computed deterministically from the markup:

\begin{itemize}[noitemsep]
  \item the page kind, from a closed vocabulary (home, list, detail, form, review, outcome,
        info, login, error, other);
  \item the URL path with identifiers replaced (\code{/orders/\{n\}});
  \item the main heading;
  \item whether the form on the page is empty, partly filled or full;
  \item whether an error or alert is shown;
  \item whether the page offers a submit action.
\end{itemize}

Controls get roles from markup cues: submit, edit, remove, add, back, navigate. Two rules
came from development evidence. Links and buttons in the page's navigation, header or
footer are never "submit": they start a new journey rather than complete the current one.
And role cues match word prefixes, so "Saving" and "Save" share a role.

Experiment 2 showed that the network does not learn these distinctions by itself, and the
literature finds that the choice of state abstraction drives results in both web
testing~\cite{chang2023webqt, liu2026webgui} and count-based
exploration~\cite{tang2017exploration}. The granularity was tuned on development
applications only: exact field combinations gave 150--225 cells on one application, so form
state was reduced to empty, partial and full, which gave about 20.

\section{Archive with exact return}
\label{sec:archive}

The tester keeps an archive of every cell it has reached, with the shortest known action
sequence to it, its visit count, and the operators already tried there. At the start of an
excursion it chooses a cell with probability proportional to a weight
\begin{equation}
  w(c) = \frac{1}{\sqrt{1 + n_c}}\cdot\Big(1 + \frac{\min(u_c, 10)}{5}\Big)\cdot
         \frac{1}{1 + 2 f_c}\cdot \rho_c \cdot \Big(1 + \beta_d\,\frac{\min(d_c, 30)}{30}\Big),
\end{equation}
where $n_c$ is the number of visits, $u_c$ the number of untried operators, $f_c$ the
number of failed returns, $\rho_c$ is 1 if enough of the episode remains after reaching the
cell and 0.1 otherwise, $d_c$ the cell's depth, and $\beta_d = 2$ a preference for depth
(faults sit many steps into a workflow). The tester then \emph{returns} to the cell by
replaying its path exactly, and acts from there. This is
Go-Explore~\cite{ecoffet2021first}: poor exploration comes from forgetting how to get back to
promising states, and the deep, unfound faults of Experiment 1 were exactly that case.

Paths are pruned of steps that change neither the page nor the cell, and every replayed
action counts towards the interaction budget, so returning is not a hidden advantage. One
replay bug, in which pruning removed steps that raised a notice the next step depended on,
was found on a development application and fixed; failed returns fell from 6 of 14 to 0 of
28.

\section{Test operators}
\label{sec:operators}

From a cell the tester applies a \emph{test operator}, a short procedure from standard test
design~\cite{myers2011art, beizer1990software}:

\begin{itemize}[noitemsep]
  \item \code{complete\_step}: fill the form validly and submit (the path deeper);
  \item \code{choose\_option}: take each option of a dropdown in turn (equivalence classes);
  \item \code{invalid\_input}: enter an empty value, a value of the wrong type, or a boundary
        minimum or maximum (boundary-value analysis);
  \item \code{edit\_after\_review}: reach a review step, go back and change an earlier
        answer (state-transition testing);
  \item \code{back\_and\_resume}, \code{reload\_and\_resume}: interrupt a flow and continue;
  \item \code{double\_submit}: submit twice in quick succession;
  \item \code{collection\_churn}: add and remove items in a collection;
  \item \code{combine} (development option, not adopted): two deviations on one form, such
        as a dropdown option together with an invalid field (pairwise testing);
  \item \code{browse\_list} (option, Section~\ref{sec:v3-options}): sort or filter a list,
        page forward, and open an item on a later page.
\end{itemize}

Each operator then \emph{follows through}: it continues the task to submission, for up to
four submissions, because a deviation usually shows its effect only on a later page. When a
page offers several submit actions (for example, "Update bag" and "Go to checkout"),
operators take them in turn on repeat visits. Operators are chosen by technique, not from
any application's faults. They are how the tester leaves the happy path, which the LLM
progress reward could not do.

\section{The decision layer}
\label{sec:decision}

The decision layer chooses which cell to return to and which operator to apply on which
target. Three variants exist; each is an arm.

\paragraph{No learner (\code{hybrid\_nolearn}).} Cells by the archive weight; at a cell,
untried (operator, target) pairs first, in random order. This is the non-learning
reference.

\paragraph{Thompson-sampling bandit (\code{hybrid\_bandit}).} A Bayesian bandit learns,
online inside the application, which operator pays off on which kind of target. Its keys
are semantic and back off from specific to general:
\[
  (\text{page kind}, \text{operator}, \text{target type}) \;\rightarrow\;
  (\text{operator}, \text{target type}) \;\rightarrow\; (\text{operator}).
\]
Each key has a $\mathrm{Beta}(\alpha, \beta)$ posterior over the probability that the
operator produces an interesting outcome. A specific key's prior is the more general keys'
evidence, discounted by 0.3, so a new combination starts from what similar ones did. At each
decision the tester samples from every candidate's posterior and takes the best
sample~\cite{thompson1933likelihood, russo2018tutorial}. The reward is the tester's own
observation, never the seeded-fault oracle: an alert shown, new cells reached, or an
invariant candidate. A bandit was chosen over a deep network for three reasons:
Experiment 1 found no gain from bootstrapped multi-step learning; a 2{,}000-interaction run
makes only about 50--135 decisions, enough for a bandit over a few dozen interpretable keys
but far too few for a network; and keys based on page kinds, not application text, can carry
across applications.

\paragraph{Curiosity Q-learning (\code{hybrid\_curiosity}).} This variant adopts the
literature's design for the free exploration steps between operators. Following
WebExplor~\cite{zheng2021webexplor} and WebRLED~\cite{gu2025webrled}, a tabular
Q-learner over cells is rewarded on every step for curiosity:
\begin{equation}
  r = \frac{1}{\sqrt{N(c, a, c')}} + 0.5\cdot\mathbb{1}[c' \text{ not yet seen this episode}],
\end{equation}
where $N(c, a, c')$ counts the transition. It learns with $\alpha = 0.5$, $\gamma = 0.9$ and
$\epsilon = 0.1$, with untried actions initialised optimistically at
$R_{\max}/(1-\gamma)$, so that an action that changes nothing cannot outrank untried ones by
bootstrapping off itself. Being off-policy, it learns from every transition, about 2{,}000
updates per run, against about 135 outcomes for the bandit. A standalone tester built only
from these components (no operators, no bandit) was also run, as
\code{dev\_curiosity\_rl}.

\section{Options for the final tuning round}
\label{sec:v3-options}

After the held-out evaluation (Chapter~\ref{chap.7}) and the rebuild of the development
applications, three options were added, each switched on in its own tuning arm:

\begin{itemize}
  \item \textbf{\code{browse\_list}} (\code{dev\_v3\_browse}). The held-out evaluation
        exposed a blind spot: the operators work on forms, and no operator pages through a
        list. This operator applies a sort or filter control, presses "next" or "load more"
        up to twice, and opens an item on the later page.
  \item \textbf{Custom clickables} (\code{dev\_v3\_clickables}). The registry also offers
        elements made clickable by ARIA roles, click handlers or styling, such as kanban
        cards and tab strips.
  \item \textbf{Failure feedback} (\code{dev\_v3\_failfeedback}). Cells where a browser
        error followed an action get extra archive weight (up to $\times 3$), since code
        that already misbehaves is likelier to misbehave again nearby.
\end{itemize}

The curiosity variant (\code{dev\_v3\_curiosity}) is the fourth tuning arm.

\section{Development rounds}
\label{sec:devrounds}

The design was tuned on development data only: the six benchmark applications, a generated
application family, Gitea, and the eight development suite applications
(Chapter~\ref{chap.6}). Each round compared arms on the same applications and seeds, and
ended with a recorded decision. Table~\ref{tab:devrounds} lists them.

\begin{table}[H]
\centering
\small
\caption{Development rounds of the hybrid tester and their decisions.}
\label{tab:devrounds}
\begin{tabularx}{\textwidth}{@{}lX X@{}}
\toprule
Round & Tested & Decision \\
\midrule
Benchmark 1 & bandit, no learner, random (4 apps) & discarded (code changed mid-round); replay bug found and fixed \\
Benchmark 2 & bandit vs no learner, with clean builds (6 apps) & validation invariant narrowed; form fixes \\
Benchmark 3 & free-exploration variants & depth preference and weighted exploration adopted \\
Benchmark 4--5 & which submit actions a cell records & cells record \emph{whether} a submit is offered \\
Benchmark 6--7 & whole-word vs prefix matching of role cues & prefix matching adopted \\
Suite 1 & bandit vs no learner, 8 suite apps & site navigation never submits; submit actions taken in turn \\
Suite 2 & the same after the fix, 14 apps & kept; learner not ahead \\
Suite 3 & two corrected bandit rewards (novelty, surprise) & no difference: reward not the bottleneck \\
Suite 4 & pairwise operator + technique-first choice & worse; not adopted \\
Suite 5 & curiosity Q-learning, in the hybrid and standalone & neither adopted \\
\bottomrule
\end{tabularx}
\end{table}

\begin{table}[H]
\centering
\small
\caption{All arms on the same runs: suite rounds 2--5, 14 applications $\times$ 3 seeds.
Difference is the per-application mean $T$ against \code{hybrid\_nolearn}, with a bootstrap
interval over applications (negative = faster).}
\label{tab:devarms}
\begin{tabular}{@{}lrrrl@{}}
\toprule
Arm & Found a fault & Median $T$ & Mean $T$ & Difference [95\% interval] \\
\midrule
\code{hybrid\_nolearn} & 32 / 42 & 742 & 889 & reference \\
\code{hybrid\_bandit} & 30 / 42 & 335 & 840 & $-49$ $[-279, +190]$ \\
\code{dev\_bandit\_surprise} & 31 / 42 & 586 & 844 & $-46$ $[-257, +143]$ \\
\code{dev\_bandit\_novelty} & 31 / 42 & 486 & 854 & $-35$ $[-252, +190]$ \\
\code{dev\_pairwise\_bandit} & 27 / 42 & 757 & 943 & $+53$ $[-121, +227]$ \\
\code{dev\_hybrid\_curiosity} & 29 / 42 & 612 & 978 & $+89$ $[-38, +202]$ \\
\code{dev\_pairwise\_nolearn} & 26 / 40 & 1{,}069 & 1{,}046 & $+119$ $[-37, +281]$ \\
\code{dev\_curiosity\_rl} & 17 / 42 & 2{,}000 & 1{,}313 & $+423$ $[+142, +744]$ \\
\code{random} & 18 / 42 & 2{,}000 & 1{,}355 & $+465$ $[+178, +790]$ \\
\bottomrule
\end{tabular}
\end{table}

\paragraph{Is the learner helping?} The top four arms in Table~\ref{tab:devarms} are
statistically tied. \code{hybrid\_nolearn} verified a fault in the most runs, and
\code{hybrid\_bandit} was fastest when it found one. No learner variant improved on no
learner: three bandit reward signals, the pairwise operator, and curiosity Q-learning. The
standalone curiosity tester was fast on shallow faults but never found the deep ones, and
was significantly worse than the hybrid, level with random. It is a tester built from the
published components, not a reproduction of WebExplor or WebRLED; their published results
(coverage and crash-type failures over long runs) are not contradicted by it.

The design leaves the learner little to decide, for four reasons. On about 5 of the 14
applications a fault is found within the first few decisions, before learning can matter.
The faults no arm finds need two deviations in one walk (for example, express delivery
\emph{and} an empty address), which choosing one operator more cleverly cannot produce. The
option space is small: about 28 cells times 8 operators can mostly be covered in about 135
decisions, so trying untried options systematically already does most of the work. And
operators with many targets crowd out the rest: invalid input took 30\% of all decisions
because pages have many fields. This agrees with the Humanoid authors' observation that
their learned model was not the main reason for their gains~\cite{li2019humanoid}: on
functional faults, the structure (abstraction, archive with exact return, operators) does
the work.
